%% kapa: exam=1 round=24 year=2013 period=1 form=A subject=civil_law subject_ko=민법 seq=23 number=23
%% source: Q-Net 감정평가사 기출문제 (cst00309) artlSeq=5073802 file=extracted/1008/5073802/2013년 제24회 감정평가사 1차시험 기출문제/1교시(A형).hwp
%% conversion: kordoc 4.24.51 -> tools/gen1.py (auto); review=proofread-sonnet; flags=none
\begin{question}{23}
미등기 부동산에 관한 법률관계에 대한 설명으로 옳지 않은 것은? (다툼이 있으면 판례에 의함)
\begin{choices}
  \item 소유자 甲으로부터 미등기 건물을 매수한 乙이 소유권이전등기를 마치지 않은 상태에서 이를 다시 丙에게 전매하여 인도한 경우 甲은 丙에게 소유물반환을 청구할 수 없다.
  \item 미등기 무허가건물을 매수하였으나 소유권이전등기를 마치지 않은 매수인은 그 건물의 불법점유자에 대하여 소유물반환을 청구할 수 없다.
  \item 미등기 건물에 대한 양도담보계약상의 채권자의 지위를 승계하여 건물을 관리하고 있는 자는 건물에 대한 철거처분권을 갖지 못한다.
  \item 미등기 건물을 그 대지와 함께 매수한 사람이 그 대지에 관하여만 소유권이전등기를 넘겨받아 저당권을 설정한 후 그 저당권의 실행으로 대지가 경매되어 다른 사람의 소유로 된 경우 그 건물을 위한 법정지상권이 성립한다.
  \item 미등기 부동산의 매수인이 그 부동산을 인도받은 경우 매도인의 상속인이 행사하는 반환청구에 대하여 반환을 거절할 수 있다.
\end{choices}
\end{question}
\answer{4}
